\begin{solucion}
Por definición el cuerpo de descomposición de un polinomio \(f(x) \in \mathbb{Q}[x]\) es el cuerpo \(L\) tal que:
\begin{itemize}
    \item \(L = \mathbb{Q}(\alpha_1, \ldots, \alpha_4)\).
    \item $f(x) = c(x - \alpha_1)(x - \alpha_2)(x - \alpha_3)(x - \alpha_4) \in L$
\end{itemize}
Para \(x^4 + 2\). Este polinomio es irreducible en \(\mathbb{Q}[x]\), pues por el criterio de Eisenstein con \(p=2\), $p \mid 2$ y \(p^2 \nmid 2\), se concluye que \(x^4 + 2\) no tiene raíces racionales. Por lo que como $x^4 + 2$ es irreducible sobre \(\mathbb{Q}\) y es el polinomio mínimo de alguna raíz \(\alpha\) en el cuerpo de descomposición $L$, entonces la extensión \(\mathbb{Q}(\alpha)\) tiene grado 4 sobre \(\mathbb{Q}\).
Ahora sea $\alpha \in L$ una raíz de $x^4 + 2$, entonces:
\begin{align*}
    \alpha^4 + 2 &= 0 \\
    \alpha^4 &= -2 \\
    \alpha^4 &= 2e^{\pi i + 2\pi i k} \quad (k \in \mathbb{Z}) \\
    \alpha &= \sqrt[4]{2}e^{\pi i/4 + \pi i k/2} \quad (k = 0, 1, 2, 3).
\end{align*}
Por lo tanto, las raíces de \(x^4 + 2\) son:
\begin{align*}
    \alpha_1 &= \sqrt[4]{2}e^{\pi i/4}, \\
    \alpha_2 &= \sqrt[4]{2}e^{3\pi i/4}, \\
    \alpha_3 &= \sqrt[4]{2}e^{5\pi i/4}, \\
    \alpha_4 &= \sqrt[4]{2}e^{7\pi i/4}.
\end{align*}
donde:
\begin{align*}
    \omega_1 &= e^{\pi i/4}= e^{2\pi i/8}, \\
    \omega_2 &= e^{3\pi i/4}= e^{2\pi i(3)/8}, \\
    \omega_3 &= e^{5\pi i/4}= e^{2\pi i(5)/8}, \\
    \omega_4 &= e^{7\pi i/4}= e^{2\pi i(7)/8}.
\end{align*}
Por lo que observamos que las raíces son $\{(\omega_1, \omega_2, \omega_3, \omega_4) \} \subset \{e^{2\pi i k/8}\}_{k=0}^{7}$, estan contenidas en las raíces octavas de la unidad. De esta forma las raíces de \(x^4 + 1\) pueden ser generadas por $i$, de este modo, tomamos $L=\mathbb{Q}(\sqrt[4]{2}, i)$ el cual es el cuerpo de descomposición, dado que $\sqrt{2}/2 \in \mathbb{Q}(\sqrt[4]{2}, i)$ entonces $\sqrt[4]{2} (\sqrt{2}/2 + i \sqrt{2}/2) = \alpha_1 \in L$, luego tambien esta $\alpha_2 = \sqrt[4]{2} (\sqrt{2}/2 - i \sqrt{2}/2) \in L$, $\alpha_3 = \sqrt[4]{2} (-\sqrt{2}/2 + i \sqrt{2}/2) \in L$ y $\alpha_4 = \sqrt[4]{2} (-\sqrt{2}/2 - i \sqrt{2}/2) \in L$. Por lo que $L$ es el cuerpo de descomposición de \(x^4 + 2\).Po r lo que usando el teorema de la torre:
\begin{align*}
    [L:\mathbb{Q}] &= [\mathbb{Q}(\sqrt[4]{2}, i):\mathbb{Q(\sqrt[4]{2})}][\mathbb{Q}(\sqrt[4]{2}):\mathbb{Q}] \\
\end{align*}
donde sabemos que \([Q(\sqrt[4]{2}):\mathbb{Q}]=4\). Asi analicemos el grado de la extensión [\(\mathbb{Q}(\sqrt[4]{2}, i):\mathbb{Q(\sqrt[4]{2})}\)] como $\mathbb Q(\sqrt[4]{2})$  no contiene raíces complejas y $x^2 + 1$ es irreducible sobre \(\mathbb{Q}(\sqrt[4]{2})\) cuya raíz es \(i\), entonces \([Q(\sqrt[4]{2}, i):\mathbb{Q}(\sqrt[4]{2})]=2\). Por lo que:
\begin{align*}
    [L:\mathbb{Q}] &= [\mathbb{Q}(\sqrt[4]{2}, i):\mathbb{Q(\sqrt[4]{2})}][\mathbb{Q}(\sqrt[4]{2}):\mathbb{Q}] \\
    &= 2 \cdot 4 = 8.
\end{align*}
\end{solucion}
